\chapter{Dental Fillings}

\section*{Overview}
Dental fillings restore a tooth damaged by decay by removing the cavity and filling it with a durable material.

\section*{Alternative Names}
Dental restoration; tooth filling; amalgam or composite filling

\section*{Principle}
Carious tissue is removed, the cavity is shaped, and it is filled with amalgam, composite resin, or other material to restore function and prevent further decay.

\section*{History}
Tooth filling has been practiced for centuries; amalgam appeared in the 19th century and modern tooth-colored composites in the late 20th century.

\section*{Why It Is Done}
Performed for dental caries, broken or worn teeth, and to restore tooth structure.

\section*{Is This a Primary or Secondary Test or Procedure}
Primary restorative treatment for dental decay.

\section*{How It Is Performed}
Performed under local anesthesia. The decay is drilled away, the cavity is cleaned and shaped, and the filling material is placed, shaped, and cured.

\section*{Preparation}
Clinical examination and radiographs to assess decay extent.

\section*{Contraindications and Precautions}
Very large cavities may require crowns or root canal treatment.

\section*{Risks}
Recurrent decay, filling failure or loss, sensitivity, and allergic reactions (rare).

\section*{Aftercare and Recovery}
The patient avoids chewing on the filling briefly; sensitivity settles.

\section*{Outcome and Follow-up}
Success is judged by a comfortable bite and no recurrent decay at review.

\section*{Expected Results}
A well-sealed, functional restoration.

\section*{Follow-up}
Routine examinations and radiographs.

\section*{When to Seek Medical Care}
Follow the procedure team's specific instructions. Seek urgent medical assessment for severe or worsening pain, uncontrolled bleeding, fever or signs of infection, trouble breathing, chest pain, fainting, new weakness or confusion, inability to urinate, or any rapidly worsening symptom. The appropriate warning signs depend on the procedure and the patient's medical condition.

\section*{References}
\begin{enumerate}
  \item MedlinePlus. Dental filling. U.S. National Library of Medicine; 2025.
  \item Current Medical Diagnosis and Treatment. McGraw-Hill; 2025.
\end{enumerate}

\clearpage
